% =====================================================
% ACADEMIC WHITEPAPER TEMPLATE
% Minimal starting point for publication-grade papers
% =====================================================

\documentclass[11pt,letterpaper]{article}

% -------------------- Preamble --------------------
% See references/latex-preamble.tex for complete package list
\usepackage[utf8]{inputenc}
\usepackage[T1]{fontenc}
\usepackage[margin=1in]{geometry}
\usepackage{microtype}

% Math
\usepackage{amsmath,amssymb,amsthm}
\usepackage{mathtools}

% Graphics
\usepackage{graphicx}
\usepackage{tikz}
\usepackage{svg}

% Code
\usepackage{listings}
\usepackage{xcolor}

% Tables
\usepackage{booktabs}

% Citations
\usepackage{cite}
\usepackage{hyperref}
\hypersetup{
    colorlinks=true,
    linkcolor=blue,
    citecolor=blue,
    urlcolor=blue
}

% Theorem environments
\newtheorem{theorem}{Theorem}[section]
\newtheorem{lemma}[theorem]{Lemma}
\newtheorem{definition}{Definition}[section]

% -------------------- Document Begin --------------------
\begin{document}

% Title Page
\begin{titlepage}
\centering

{\LARGE\bfseries Your Whitepaper Title}\\[0.5em]
{\Large Optional Subtitle}\\[2em]

{\large Your Name}\\
{\normalsize Your Affiliation}\\[1em]

{\normalsize Date: \today}\\
{\normalsize Version: v1.0}\\[2em]

\vfill

\begin{abstract}
Your abstract goes here. Summarize the key contributions, methodology, 
and findings in 150-250 words. This section should be self-contained 
and enable readers to determine paper relevance without reading further.
\end{abstract}

\end{titlepage}

% Table of Contents
\tableofcontents
\newpage

% -------------------- Sections --------------------

\section{Introduction}

Establish the problem context, motivate the research question, and preview 
the paper's structure. Connect to existing literature and identify gaps 
that this work addresses.

\subsection{Background}

Provide essential background information necessary for understanding the 
contribution.

\subsection{Contributions}

Enumerate the specific contributions of this work:

\begin{enumerate}
    \item First major contribution
    \item Second major contribution
    \item Third major contribution
\end{enumerate}

\section{Mathematical Foundations}

Present the formal mathematical framework.

\begin{definition}[Fixed Point]
Let $f: X \to X$ be a function on a set $X$. A point $x^* \in X$ 
is a \emph{fixed point} of $f$ if $f(x^*) = x^*$.
\end{definition}

\begin{theorem}[Banach Fixed Point Theorem]
Let $(X, d)$ be a complete metric space and $f: X \to X$ a contraction 
mapping with constant $0 \le k < 1$. Then $f$ has a unique fixed point 
$x^* \in X$.
\end{theorem}

\begin{proof}
The proof follows by iterative application of the contraction property...
\end{proof}

\section{Methodology}

Describe the technical approach, algorithms, and experimental design.

\subsection{Algorithm Description}

\begin{algorithm}
\caption{Fixed Point Iteration}
\begin{algorithmic}[1]
\Procedure{FixedPoint}{$f, x_0, \epsilon$}
    \State $x \gets x_0$
    \While{$d(f(x), x) > \epsilon$}
        \State $x \gets f(x)$
    \EndWhile
    \State \Return $x$
\EndProcedure
\end{algorithmic}
\end{algorithm}

\section{Implementation}

Reference executable code examples:

\lstinputlisting[language=Python, caption={Fixed Point Implementation}]{../src/code-examples/fixed_point.py}

\section{Results}

Present empirical findings with visual support.

\begin{figure}[htbp]
\centering
\includesvg[width=0.8\textwidth]{../src/visuals/convergence-plot}
\caption{Convergence rate showing exponential decay of distance metric 
over iterations.}
\label{fig:convergence}
\end{figure}

As shown in Figure~\ref{fig:convergence}, the algorithm exhibits 
exponential convergence...

\section{Discussion}

Interpret results, discuss limitations, and place findings in broader context.

\subsection{Comparison with Prior Work}

Compare with existing approaches, highlighting advantages and tradeoffs.

\subsection{Limitations}

Acknowledge limitations of the current approach and potential extensions.

\section{Conclusion}

Summarize key findings and outline future research directions.

% -------------------- Bibliography --------------------
\bibliographystyle{ieeetr}
\bibliography{references}

% -------------------- Appendices --------------------
\appendix

\section{Additional Proofs}

Supplementary mathematical proofs and derivations.

\section{Code Listings}

Complete code listings for reproducibility.

\end{document}
